\documentclass{article}
\usepackage[utf8]{inputenc}
\usepackage[T1]{fontenc}
\usepackage{geometry}
\usepackage{graphicx}
\usepackage{float}
\usepackage{hyperref}
\usepackage[english,indonesian]{babel}

\title{Optimasi Pemrosesan Template LaTeX Menggunakan Pendekatan Berbasis AI}
\author{Academic Support Agent}
\date{\today}

\begin{document}
\maketitle

\begin{abstract}

\end{abstract}

\chapter{PENDAHULUAN}


\chapter{TINJAUAN PUSTAKA}


% Original content (commented out):
% \documentclass[11pt,letterpaper]{article}
\usepackage[utf8]{inputenc}
\usepackage[spanish, es-tabla]{babel}
\usepackage[version=3]{mhchem}
\usepackage[journal=jacs]{chemstyle}
\usepackage{amsmath}
\usepackage{amsfonts}
\usepackage{amssymb}
\usepackage{makeidx}
\usepackage{xcolor}
\usepackage[stable]{footmisc}
\usepackage[section]{placeins}
%Paquetes necesarios para tablas
\usepackage{longtable}
\usepackage{array}
\usepackage{xtab}
\usepackage{multirow}
\usepackage{colortab}
%Paquete para el manejo de las unidades
\usepackage{siunitx}
\sisetup{mode=text, output-decimal-marker = {,}, per-mode = symbol, qualifier-mode = phrase, qualifier-phrase = { de }, list-units = brackets, range-units = brackets, range-phrase = --}
\DeclareSIUnit[number-unit-product = \;] \atmosphere{atm}
\DeclareSIUnit[number-unit-product = \;] \pound{lb}
\DeclareSIUnit[number-unit-product = \;] \inch{"}
\DeclareSIUnit[number-unit-product = \;] \foot{ft}
\DeclareSIUnit[number-unit-product = \;] \yard{yd}
\DeclareSIUnit[number-unit-product = \;] \mile{mi}
\DeclareSIUnit[number-unit-product = \;] \pint{pt}
\DeclareSIUnit[number-unit-product = \;] \quart{qt}
\DeclareSIUnit[number-unit-product = \;] \flounce{fl-oz}
\DeclareSIUnit[number-unit-product = \;] \ounce{oz}
\DeclareSIUnit[number-unit-product = \;] \degreeFahrenheit{\SIUnitSymbolDegree F}
\DeclareSIUnit[number-unit-product = \;] \degreeRankine{\SIUnitSymbolDegree R}
\DeclareSIUnit[number-unit-product = \;] \usgallon{galón}
\DeclareSIUnit[number-unit-product = \;] \uma{uma}
\DeclareSIUnit[number-unit-product = \;] \ppm{ppm}
\DeclareSIUnit[number-unit-product = \;] \eqg{eq-g}
\DeclareSIUnit[number-unit-product = \;] \normal{\eqg\per\liter\of{solución}}
\DeclareSIUnit[number-unit-product = \;] \molal{\mole\per\kilo\gram\of{solvente}}
\usepackage{cancel}
%Paquetes necesarios para imágenes, pies de página, etc.
\usepackage{graphicx}
\usepackage{lmodern}
\usepackage{fancyhdr}
\usepackage[left=3.17cm,right=3.17cm,top=2.54cm,bottom=2.54cm]{geometry}

%Instrucción para evitar la indentación
\setlength\parindent{0pt}
%Paquete para incluir la bibliografía
\usepackage[backend=bibtex,style=chem-acs,biblabel=dot]{biblatex}
\addbibresource{Mendeley.bib}

%Tipo de fuente arial
\renewcommand{\rmdefault}{phv} % Arial
\renewcommand{\sfdefault}{phv} % Arial

%Formato del título de las secciones

\usepackage{titlesec}
\usepackage{enumitem}
\titleformat*{\section}{\bfseries\large}
\titleformat*{\subsection}{\bfseries\normalsize}

%Creación del ambiente anexos
\usepackage{float}
\floatstyle{plaintop}
\newfloat{anexo}{thp}{anx}
\floatname{anexo}{Anexo}
\restylefloat{anexo}
\restylefloat{figure}

%Modificación del formato de los captions
\usepackage[margin=10pt,labelfont=bf]{caption}

%Paquete para incluir comentarios
\usepackage{todonotes}

%Paquete para incluir hipervínculos
\usepackage[colorlinks=true, 
            linkcolor = black,
            urlcolor  = black,
            citecolor = black,
            anchorcolor = blue]{hyperref}

%%%%%%%%%%%%%%%%%%%%%%
%Inicio del documento%
%%%%%%%%%%%%%%%%%%%%%%

\begin{document}
\pagestyle{empty}

\renewcommand{\labelitemi}{$\checkmark$}

\renewcommand{\CancelColor}{\color{red}}

\newcolumntype{L}[1]{>{\raggedright\let\newline\\\arraybackslash}m{#1}}

\newcolumntype{C}[1]{>{\centering\let\newline\\\arraybackslash}m{#1}}

\newcolumntype{R}[1]{>{\raggedleft\let\newline\\\arraybackslash}m{#1}}

\begin{center}

%Reemplace el texto en cada uno de los siguientes espacios. No borre los corchetes porque se producen errores en la compilación del documento.

\textbf{\LARGE{ESPECIALIZACIÓN EN QUÍMICA AMBIENTAL}}\\
	%\vspace{3mm}
	\textbf{\large{ESCUELA DE QUÍMICA}}\\
	\textbf{\large{UNIVERSIDAD INDUSTRIAL DE SANTANDER}}\\
	\vspace{8mm}
	\textbf{\normalsize{FORMATO DE PRESENTACIÓN DEL PLAN DE MONOGRAFÍA}}\\
\end{center}
\vspace{4mm}

%%%%%%%%%%%%%%%%%%%%%%%%%%%%%%%%%%%%%%%%%%%%%%%%%%%%%%%%%%%%%%%%%%%%%%%%%%%%%%%%%%%%%%%%%%%%%%%%%%%%%%%%%%%%%%%%%%%%%%%%%%%%%%%%%%%%%%%%%Inicio del documento

TÍTULO DE LA MONOGRAFÍA: UN TÍTULO MUY INTERESANTE \vspace{4mm} \\ 
NOMBRE DEL ESTUDIANTE: PEPITO PÉREZ \vspace{4mm} \\
FECHA DE ENTREGA: 12 de marzo de 2018 \vspace{4mm} \\
DIRECTOR: JOSÉ PEKERMAN \vspace{4mm} \\
FIRMA DEL DIRECTOR: 
\vspace{4mm}

\section{RESUMEN DE LA PROPUESTA (Máximo 100 palabras)} 

Descripción de la estructura, objetivos e impacto del proyecto.

\section{INTRODUCCIÓN (Máximo 350 palabras)}

Marco de referencia para establecer la importancia de la ejecución de la propuesta.

\section{PLANTEAMIENTO DEL PROBLEMA (Máximo 250 palabras)}

Exposición clara, concisa y bien fundamentada del problema que se plantea. Se usan algunas referencias generales  (soportes documentales, estadísticas, datos, bibliografía, etc.) para establecer la importancia y la pertinencia del tema de trabajo propuesto. Para colocar las referencias el primer paso es vincular la cuenta de Mendeley con la cuenta de Overleaf y para citar se utiliza el comando: \textit{autocite} La referencia aparece automáticamente como un superíndice en el punto de inserción y el texto completo aparece al final en la sección de referencias. Ejemplo de referencia \autocite{young:1902}. Otro ejemplo \autocite{Lozano2010}

%Ejemplo de tabla

\begin{table}
\caption{Título de la tabla} %Título de la tabla
\label{tab:ejemplo} %Etiqueta para referenciar la tabla
\begin{tabular}{L{4cm}C{4cm}R{4cm}} %Definición de las columnas C{4cm} significa una columna de 3cm de ancho y con el texto centrado. L{4cm} y R{4cm} es alineación a la izquierda y a la derecha respectivamente.
\hline % Este comando dibuja una línea horizontal
\textbf{Encabezado 1} & \textbf{Encabezado 2} & \textbf{Encabezado 3} \\ %El comando \textbf sirve para escribir texto en negrilla.
\hline
1 & 2 & 3 \\ %Cada línea de la tabla termina con un doble backslash
4 & 5 & 6 \\
7 & 8 & 9 \\
\hline
\end{tabular}
\end{table}

Las tablas se pueden referenciar dentro del texto utilizando el comando:
\begin{verbatim}
\ref{etiqueta de la tabla}
\end{verbatim}

Como se observa en la \ref{tab:ejemplo} existe una relación lineal entre las diferentes columnas.

\section{MARCO TEÓRICO (Máximo 750 palabras)}

Fundamentación de la propuesta basada en recopilación y análisis de la información bibliográfica.  La revisión del estado del arte en esta etapa se hace utilizando bases de datos, libros y artículos científicos. En esta sección de la propuesta son de vital importancia las citaciones y las referencias bibliográficas para soportar las afirmaciones hechas en el documento. \\

Para incluir imágenes dentro del texto debemos importar los archivos de en formato .jpeg sin son fotos o en formato vectorial (.eps, pdf, etc.) las ilustraciones, diagramas, etc. 

\begin{figure}
\caption{Título de la imagen}
\label{fig:loro}
\includegraphics[width=15cm]{loro.jpeg}\\ %El tamaño de la imagen se puede controlar con la instrucción [width=10cm] que significa que la imagen será insertada dentro del texto con un ancho de 10 cm.
\raggedright
\hspace{10pt} \small{FUENTE: Descargado de Internet bajo licencia \textit{Creative Commons}}
\end{figure}

Al igual que con las tablas, las imágenes también las podemos referenciar utilizando el comando \textit{ref}. Por ejemplo para citar la fotografía de los loros utilizamos el comando \ref{fig:loro}.

\section{OBJETIVOS}

\subsection{Objetivo General}

El objetivo general resume la idea completa de la propuesta.  Casi siempre está muy relacionado con el título de la propuesta. La redacción de los objetivos específicos es muy importante porque establece los parámetros sobre los cuales se realiza la evaluación de la propuesta por parte de los jurados. Consecuentemente deben cumplir con los siguientes criterios:

\subsection{Objetivos Específicos}

\begin{itemize}
\item Específicos. Qué, cómo, cuándo, por qué?
\item Medibles. Cuánto?
\item Alcanzables. Habilidades, esfuerzos y limitaciones
\item Realistas. Se pueden lograr. 
\item Tangible.  Se puede apreciar.
\end{itemize}

\section{METODOLOGÍA (Máximo 750 palabras)}

En esta sección se presenta la aproximación metodológica a la solución del problema.  Se detallan los recursos técnico-científicos y las actividades necesarias para lograr los objetivos propuestos.  Se pueden utilizar apéndices, diagramas, para aclarar secuencias.

\section{RESULTADOS ESPERADOS (Máximo 100 palabras)}

Listado de los posibles resultados que se pretende alcanzar con la 	ejecución de la propuesta.

\section{CRONOGRAMA}

Tabla conteniendo el tiempo necesario para realizar la propuesta que incluya un desglose detallado por objetivos específicos o por actividades. 

\section{PRESUPUESTO}

Desglose y justificación de los gastos necesarios para realizar los objetivos de la propuesta.

\section{REFERENCIAS}

%No borrar este comando porque desaparece la bibliografía.
\printbibliography[heading=none]
		
\end{document}


\end{document}